\section{Results and Discussion}
\subsection{Preliminary Measured-data Evidence}
The separate non-CAMS run contains 24 temporal evaluations (expanding and
sliding sensitivities), 36 spatial holdouts and 12 calibration evaluations.
Table~\ref{tab:paired} and Fig.~\ref{fig:paired} show two paired development
comparisons. Strict Ridge improves climatology on these pairs; history Ridge
does not establish improvement over persistence. These comparisons answer
neither the raw-CAMS correction question nor the locked final-test question.
\begin{table}[t]\centering\footnotesize
\caption{Non-CAMS preliminary paired MAE differences. Negative differences favour Ridge; intervals use seven-day date blocks.}
\label{tab:paired}
\setlength{\tabcolsep}{4pt}
\begin{tabular}{lrrr}
\toprule
Comparison & $\Delta$MAE & 95\% interval & $n$\\\midrule
Strict Ridge--climatology & -8.96 & [-12.49, -5.61] & 38,382\\
History Ridge--persistence & +0.61 & [-0.41, +1.70] & 31,622\\
\bottomrule\end{tabular}
\end{table}
\begin{figure}[t]\centering\AirGuardPaired
\caption{Paired non-CAMS development differences. The persistence comparison
crosses zero; a lower error than climatology does not establish a CAMS improvement.}
\label{fig:paired}\end{figure}

\subsection{Uncertainty and Extreme-event Failures}
First-fold marginal interval coverage masks a high-pollution failure.
Both nominal intervals cover none of the 35 exact Very Poor outcomes or
48 exact Poor-or-worse pairs. The individual Poor subgroup has only 13 pairs
and its subgroup metric is withheld. Figure~\ref{fig:coverage} retains these
denominators rather than treating unavailable estimates as zeros.
\begin{table}[t]\centering\footnotesize
\caption{First-fold non-CAMS interval diagnostics. Width is in AQI units; neither row establishes operational calibration.}
\label{tab:coverage}
\setlength{\tabcolsep}{4pt}
\begin{tabular}{lrrr}
\toprule
Nominal & Exact coverage & Mean width & $n$\\\midrule
80\% & 0.855 & 95.5 & 2,923\\
90\% & 0.944 & 120.6 & 2,923\\
\bottomrule\end{tabular}
\end{table}
\begin{figure}[t]\centering\AirGuardCoverage
\caption{Non-CAMS first-fold coverage by observed category. Both Very Poor
estimates are zero; the Poor metric is unavailable at $n<30$. Satisfactory and
Moderate are abbreviated on the axis.}\label{fig:coverage}\end{figure}

\subsection{Primary CAMS Status}
The request manifest verifies 6 of 36 chunks. The guarded join covers 42 dates and 7,560 station-days, retaining 2,054 rejected targets and 2 censored Severe targets. This coverage cannot support the locked training and verification
windows. No primary model has been fitted or selected. Missing archive months
are not replaced with reanalysis, later operational fields or fabricated data.
\textit{Primary CAMS comparison: [RESULT PENDING CAMS ARCHIVE EXPERIMENT].}

\subsection{Prospective Application}
The audited ledger has 24 immutable forecasts across 2 target days and 0 verified days. 0 forecasts have pre-issue station bindings; 12 historical entries lack an issuance Git SHA. Missing historical identities or station bindings are not inferred
after the outcome. The production application displays raw CAMS-based estimates;
no experimental correction or unvalidated interval is deployed. Its interface
and API checks establish software behaviour, while predictive skill requires
matched outcomes and the declared research evaluation.

\subsection{Calibration Stability Across Folds}
Figure~\ref{fig:width} displays mean interval widths for the existing non-CAMS
folds. These are fold-level summaries, not a reconstructed histogram of
individual intervals. Both width and coverage must be inspected: wider
intervals can improve marginal coverage while retaining poor subgroup
performance. A censored Severe outcome identifies a lower bound, so whether
an upper endpoint reaches 401 is reported separately from exact coverage.
\begin{figure}[t]\centering\AirGuardWidths
\caption{Non-CAMS preliminary mean widths across chronological folds.
The calibrated intervals remain development diagnostics rather than
validated operational intervals.}\label{fig:width}\end{figure}

\subsection{Evidence Boundaries}
The initial 90-day training fold cannot meet a station's 120-valid-day
eligibility requirement. It is preserved as insufficient rather than
rescued using later coverage. The preliminary non-CAMS study's separate
30-day station threshold and fixed penalty therefore must not be conflated
with primary tuning. Likewise, paired regression evidence is restricted to
exact targets; censored observations are retained for the declared category
and lower-bound analysis.

The comparison with persistence is an important negative finding.
A confidence interval containing zero does not demonstrate that two methods
are identical, but it also does not support claiming a reliable difference.
The first-fold Very Poor interval failure demonstrates why a single average
coverage number cannot justify high-pollution warnings. These observations
support the evaluation design, not a claim that station correction already
improves genuine archived CAMS.
